\section{Introduction}\label{sec:introduction}

A phase transition is a singularity that appears in the thermodynamic
limit of a particle system.  We construct one for a radial pair
interaction with a bounded continuous core and an explicit integrable
power-law bound on its tail.

We consider identical classical particles in three-dimensional Euclidean space.
For a radial pair potential \(\phi\), a cube \(\Lambda_L=[0,L]^3\), and \(N\ge0\), set
\begin{equation}\label{eq:partition}
 U^\phi(x_1,\ldots,x_N)=\sum_{1\le i<j\le N}\phi(\abs{x_i-x_j}),
 \qquad
 Z^\phi_{L,N}(\beta)=\frac1{N!}
  \int_{\Lambda_L^N}\exp(-\beta U^\phi)\,\dd x_1\cdots\dd x_N,
\end{equation}
with \(Z^\phi_{L,0}=1\).
The position measure is fixed independently of the inverse temperature
\(\beta>0\).
A potential is \emph{stable} if \(U^\phi\ge-BN\) for a constant \(B\)
and every finite configuration.

Let \(p\) denote the packing constant
\[
 p=\lim_{L\to\infty}\frac{M_L}{L^3},
 \qquad
 M_L=\max\bigl\{\#X:X\subset[0,L]^3,\
               \abs{x-y}\ge1\ \text{for distinct }x,y\in X\bigr\}.
\]
The existence and positivity of \(p\) are proved in \Cref{f:packing};
its numerical value is not needed.

\begin{theorem}\label{thm:main}
There is a bounded continuous function
\(\phi:[0,\infty)\to\R\) and finite constants \(B,C>0\) such that
\[
 \sum_{i<j}\phi(\abs{x_i-x_j})\ge-BN,
 \qquad
 \abs{\phi(r)}\le C r^{-3-1/32}\quad(r\ge1).
\]
There are an open interval \(I\subset(0,\infty)\) centered at \(5p/3\)
and a common inverse temperature \(\beta_c\in[7/8,9/8]\) such that,
for every \(\rho\in I\), the canonical free energy
\[
 f(\beta,\rho)=
 -\lim_{L\to\infty}\frac1{\beta L^3}
       \log Z^\phi_{L,\lfloor\rho L^3\rfloor}(\beta)
\]
exists and is finite for every \(\beta>0\).
Its finite one-sided derivatives at \(\beta_c\) satisfy
\[
 f'_-(\beta_c,\rho)>f'_+(\beta_c,\rho).
\]
In particular, \(f(\,\cdot\,,\rho)\) is not differentiable at \(\beta_c\).
\end{theorem}

\subsection{Context and the role of the tail}

The existence of continuum phase transitions belongs to the program
highlighted by Simon's problems in mathematical physics
\citep{Simon1984}.  We address the explicit interaction class and
canonical singularity in \Cref{thm:main}.

Rigorous continuum coexistence results precede the present construction.
Ruelle proved a transition for a symmetric two-species classical system
\citep{Ruelle1971}.  Lebowitz, Mazel, and Presutti established a
liquid--vapor transition for a continuum model with an attractive
pair term and a repulsive four-body term
\citep{LebowitzMazelPresutti1998,LebowitzMazelPresutti1999}.
Dereudre and Renaud-Chan establish density coexistence for a class of
saturated continuum interactions
\citep[Theorems~1--2]{DereudreRenaudChan2026}.
Their radial-kernel example integrates over the union of particle
neighborhoods, giving a many-body energy
\citep[Section~2.3, Example~3]{DereudreRenaudChan2026}.
The species or many-body structures are part of these models; the
problem considered here uses a single species and a sum of radial pair
interactions.

Long-range weak attractions provide a classical route from microscopic
particle interactions to a nonconvex mean-field free energy.  The
Lebowitz--Penrose theory identifies the limiting free energy through
convexification of a reference free energy with a quadratic attraction
term \citep{LebowitzPenrose1966}.  Gerardi, Marchioro, Olivieri, and
Presutti extended this framework from hard-core references to
superstable interactions \citep{GerardiEtAl1973}.  In the present problem, the limit of
an interaction as its range tends to infinity is only an intermediate
object: the required singularity must belong to one fixed pair
potential.  Moreover, a pointwise power bound constrains how the masses
of its successive attractions may depend on their ranges.

The companion article \citep[Theorem~1.1]{OpenAICompanion2026} constructs
a radial potential with a divergent repulsive core, an integrable
\(o(r^{-3})\) tail, and a common transition inverse temperature for an
open interval of densities.  Both articles use small clusters and
separated supports of a pressure in two variables.  Here we prove the
bounded-core construction and the quantitative range estimate needed
for a fixed power margin.  All proofs required for the present theorem
are included below.  The two articles concern different potentials,
each with its own density interval and common critical inverse temperature.

A recent theorem of He, Jauslin, Lebowitz, and Peled proves a
liquid--vapor transition for a finite-range modified Kac pair
interaction \citep[Equation~(1.13), Theorems~1.3 and~1.5]{HeJauslinLebowitzPeled2026}.
Their interaction is defined using a fixed partition into boxes, and
the coexistence theorem is formulated at a fixed inverse temperature.
Our physical interaction is a radial function of particle separation;
boxes enter only the estimates used to analyze its integrals.  The
conclusion here is a strict inverse-temperature derivative jump of the
canonical free energy at one common critical point throughout an open
interval of densities.  It is an existence result for
the constructed infinite-range interaction, with no identification of
its phases as conventional fluid phases.

The estimates also have several classical ancestors.  The quadratic
occupation control is part of the superstability framework developed by Ruelle \citep{Ruelle1970}.
For our particular core, the needed bound follows from the weighted
graph maximization argument of Motzkin and Straus
\citep[Theorem~1]{MotzkinStraus1965}; we reproduce the argument and its
geometric application.  Supporting slopes and partial convex duality
are used in their standard finite-dimensional sense
\citep[Theorem~11.1 and Proposition~11.3]{RockafellarWets1998}.
The local occupation estimates, the quantitative attraction comparison,
and the persistence argument are proved at their precise strengths in
this paper.

\subsection{Proof strategy}

There are three tasks: construct a stable family of bounded potentials,
compare a long-range attraction with its ideal density transform, and
preserve a sufficiently large set of pressure supports while the
attraction ranges grow.  We describe how their outputs fit together.

\paragraph{A bounded reference family.}
A parameter \(t>1\) fixes the microscopic core; it is chosen once and is
independent of \(\beta\).  Below distance one, the core has height
\(t^2\) outside a narrow shell at distances comparable to \(\ee^{-t}\);
inside the shell it has a zero-energy plateau.  It vanishes beyond a
fixed radius.  Four particles can form a tetrahedron with all six pair
distances in the zero plateau, whereas five cannot lie in the shell.  This gives a quadratic lower bound
in the particle counts of small cubes.  The bound survives the addition
of nonnegative radial attractions with a prescribed total mass.
\Cref{f:section,f:local-section} use it to establish the canonical and
grand-canonical limits, uniform occupation moments, and continuity of
pressure under a small change of attraction mass.

For a potential \(\varphi\) in this family and a density \(d\ge0\),
write
\[
 H_\varphi(\beta,d)
 =-\lim_{L\to\infty}\frac{1}{tL^3}
        \log Z^\varphi_{L,\lfloor dL^3\rfloor}(\beta),
 \qquad
 Q_\varphi(\beta,h)
 =\sup_{d\ge0}\{hd-H_\varphi(\beta,d)\}.
\]
Thus \(H_\varphi=\beta f_\varphi/t\), and \(h\) is the field dual to
density.  The pressure \(Q_\varphi\) is convex jointly in
\((\beta,h)\).  A supporting slope \((e,d)\) at \((\beta_0,h_0)\)
means that
\[
 Q_\varphi(\beta,h)\ge Q_\varphi(\beta_0,h_0)
             +e(\beta-\beta_0)+d(h-h_0)
\]
for every \(\beta>0\) and \(h\in\R\).
Two such slopes with the same density \(d=\rho\) and distinct first
coordinates give two distinct upper supporting slopes of
\(H_\varphi(\,\cdot\,,\rho)\).  This is the route to the derivative
jump, proved explicitly in \Cref{m:canonical-corner}.

\paragraph{An attraction with controlled error.}
For a pressure \(Q\), an ideal attraction of mass \(a>0\) acts by
\begin{equation}\label{eq:intro-transform}
 (T_aQ)(\beta,h)
 =\sup_{u\in\R}
   \left\{Q(\beta,h+\beta au)-\frac{\beta a}{2}u^2\right\}.
\end{equation}
Square completion identifies its maximizing shifts \(u\) with
maximizing densities after addition of the quadratic attraction term.
\Cref{k:section-transform} proves attainment, uniform bounds on these
shifts, and an exact formula for all supporting slopes of \(T_aQ\).

Fix a nonnegative radial \(C^1\) kernel \(w\), supported in the unit
ball and of integral one, and put \(w_R(x)=R^{-3}w(x/R)\).
If the reference interaction \(\varphi_0\) has range at most \(R_0\),
then \Cref{k:range} compares the physical attraction \(-ta w_R\)
with the ideal transform:
\[
 \bigl\|Q_{\varphi_0-ta w_R}-T_aQ_{\varphi_0}\bigr\|_{\mathcal D}
 \le C_{\mathcal D,t}\sqrt{R_0/R}
\]
on each compact parameter set \(\mathcal D\).  The constant is
uniform over earlier reference interactions and over positive masses
\(a\) tending to zero.  Boxes of side comparable to
\(\sqrt{R_0R}\) balance the cost of deleting reference interactions
across boxes against the variation of the new kernel inside them.
For the lower bound, independent reference grand ensembles are placed
in these boxes.  Completing a square avoids any need to identify their
mean density with the maximizing shift, so the comparison remains valid
at coexistence without selecting a phase.

\paragraph{Three slope groups and one fixed potential.}
As \(t\) grows, the bare-core pressure tends to
\(\max\{0,ph,4ph-9p\}\).  Its branches come from vacuum, single
particles, and four-particle clusters, respectively.  A first ideal
attraction turns these into three affine branches with noncollinear
slopes in the \((e,d)\) plane.  \Cref{i:section} proves that, for one
large finite \(t\), the actual generating supports remain near these
three slopes throughout a parameter neighborhood.

Uniform convergence of convex pressures by itself need not preserve
a corner.  The additional control comes from the following property
of a maximizing shift in \eqref{eq:intro-transform}: at its shifted
parameter \((\beta,h+\beta au)\), the reference pressure has a
parabolic upper bound in the field variable
(\Cref{k:active-parabola}).  In \Cref{m:one-step} this bound confines
nearby supports to one density band and hence to one of the separated
slope groups.  An error \(\delta\) in pressure then costs
\(O(\delta/a)\) in the slopes.

After the first attraction, take \(R_{n+1}=R_n^2\) and
\(a_n=R_n^{-1/16}\).  The range errors are \(O(R_n^{-1/4})\), and
their ratios to the next masses are \(O(R_n^{-1/8})\).  All relevant
losses are summable.  \Cref{sec:iteration} consequently preserves a
common two-dimensional ball of pressure supports, passes it to the
limit, and uses horizontal slices near density \(5p/3\) to obtain the
canonical singularity at a common inverse temperature throughout a
density interval.  The same schedule gives
\(\sum_n a_nR_n^{1/32}<\infty\), which supplies the pointwise tail
bound.  The core, all attraction parameters, and the density interval
are fixed before the limiting critical parameter is selected.
